% Source: https://artofproblemsolving.com/community/c3h2138952_equal_collections
% Stable identity: aops-topic-2138952-q01
\textbf{(a)} The set $B$ is the union of the disjoint sets $A$ and $B\setminus A$ (the elements of $B$ not in $A$), because $A\subseteq B$. So
\[|B|=|A|+|B\setminus A| .\]
Since $|B|=|A|$, we get $|B\setminus A|=0$: every element of $B$ is in $A$. Together with $A\subseteq B$, this gives $A=B$. $\blacksquare$

\textbf{(b)} Let $B=\{1,2,3,\dots\}$ be the positive integers and $A=\{2,4,6,\dots\}$ the even ones. Then $A\subseteq B$ and $A\ne B$ (for example $1\notin A$), but pairing each $n\in B$ with $2n\in A$ matches the elements of the two sets one-to-one. So two infinite sets can have the same size even though one is a proper part of the other.

\emph{Remark.} For collections in which elements may repeat (multisets), the statement fails even for finite collections: $\{1,1,2\}$ and $\{1,2,2\}$ have three elements each, and every element of the first appears in the second.
